% Part: counterfactuals
% Chapter: introduction
% Section: counterfactuals

\documentclass[../../../include/open-logic-section]{subfiles}

\begin{document}

\olfileid{cnt}{int}{cnt}

\olsection{प्रतिवास्तविक विधाने}

इंग्रजीतील ‘‘जर \dots तर \dots’’
या रचनेचा एक अतिशय प्रचलित आणि
महत्त्वाचा प्रकार \emph{to be}
या क्रियापदाच्या भूतकालीन
संकेतार्थी रूपाने बनतो:
‘‘जर असे घडले असते की \dots,
तर असे घडले असते की \dots’’.
अशा सशर्त विधानाचे पूर्वांग
सहसा असत्य, म्हणजे वस्तुस्थितीच्या
विरुद्ध असल्याने, त्यांना
\emph{प्रतिवास्तविक सशर्त विधाने}
म्हणतात. \emph{to be} चे
संकेतार्थी रूप त्यांत असल्याने
त्यांना \emph{संकेतार्थी सशर्त विधाने}
असेही म्हणतात. याउलट
\emph{निश्चयार्थी} सशर्त विधाने
‘‘जर असे आहे की \dots,
तर असे आहे की \dots’’
या रूपात असतात. प्रतिवास्तविक
आणि निश्चयार्थी सशर्त विधानांच्या
सत्यतेच्या अटी वेगळ्या असतात.
अॅडम्स यांचे प्रसिद्ध उदाहरण पाहा:
\begin{quote}
  जर ऑस्वाल्डने केनेडींची हत्या केली नाही,
  तर दुसऱ्या कोणीतरी केली.

  जर ऑस्वाल्डने केनेडींची हत्या केली
  नसती, तर दुसऱ्या कोणीतरी केली असती.
\end{quote}
पहिले विधान निश्चयार्थी, तर दुसरे
प्रतिवास्तविक आहे. पहिले विधान
स्पष्टपणे सत्य आहे: अध्यक्ष जॉन
एफ. केनेडी यांची हत्या
\emph{कोणीतरी} केली हे आपल्याला
माहीत आहे. ती व्यक्ती,
वॉरेन अहवालात म्हटल्याप्रमाणे,
ली हार्वे ऑस्वाल्ड नसेल,
तर दुसऱ्या कोणीतरी हत्या
केली. दुसरे विधान मात्र
वेगळे काही सांगते. त्याचा
दावा असा की ऑस्वाल्डने
केनेडींची हत्या केली नसती
म्हणजे डॅलस येथील गोळीबार
टळला असता किंवा अयशस्वी
ठरला असता, तरी पुढचा
इतिहास अशा प्रकारे घडला
असता की दुसरा हत्येचा
प्रयत्न यशस्वी झाला असता.
हे सत्य असण्यासाठी
केनेडींचा मृत्यू निश्चित
व्हावा म्हणून काही प्रबळ
शक्तींनी कट रचलेला असावा
लागेल (अनेक कटसिद्धांतवादी
जसे मानतात तसे).

\emph{निश्चयार्थी} सशर्त विधान
वास्तविक अभिव्यंजनाने अचूक
दर्शवले जाते का, हा अजून
वादाचा विषय आहे. विशेषतः,
इंग्रजी निश्चयार्थी सशर्त
विधानांच्या सत्यतेच्या अटी
वास्तविक अभिव्यंजन देते
या मताशी सुसंगत राहील
अशा प्रकारे त्याचे
विरोधाभास ‘‘समजावता’’
येतात का, हा प्रश्न
आहे. याउलट प्रतिवास्तविक
सशर्त विधाने वास्तविक
अभिव्यंजनाने अचूक
दर्शवता येत नाहीत,
यावर वाद नाही.
कारण प्रतिवास्तविक
विधानांची पूर्वांगे
सहसा असत्य असली,
तरी असत्य पूर्वांग
असलेले प्रत्येक
प्रतिवास्तविक विधान
सत्य नसते. उदाहरणार्थ,
वॉरेन अहवाल खरा
मानला आणि केनेडींच्या
हत्येसाठी कोणताही
कट नव्हता असे
मानले, तर अॅडम्स
यांचे प्रतिवास्तविक
सशर्त विधान असेच
उदाहरण ठरते.

प्रतिवास्तविक सशर्त विधाने
कार्यकारणाविषयीच्या
युक्तिवादात महत्त्वाची
असतात: प्रतिवास्तविक
विधानांनी कारणसंबंध
दर्शवणे हे त्याचे
प्रमुख उदाहरण आहे.
उदाहरणार्थ, आगकाडी
ओढल्याने ती पेटते;
हे ‘‘ही आगकाडी
ओढली असती, तर ती
पेटली असती’’ असे
म्हणून दर्शवता येते.
वास्तविक अभिव्यंजनाने,
किंवा सर्वसाधारणपणे
निश्चयार्थी सशर्त विधानांनी,
हे दर्शवता येत नाही:
‘‘आगकाडी ओढली जाते
$\lif$ आगकाडी पेटते’’
हे विधान आगकाडी
कधीच ओढली गेली
नाही तर सत्य असते,
ती ओढली असती
तर काय झाले असते
याचा काहीही संबंध
नसतो. त्याहूनही
वाईट म्हणजे,
‘‘आगकाडी ओढली जाते
$\lif$ आगकाडीचे
फुलांच्या गुच्छात
रूपांतर होते’’ हेही
आगकाडी कधीच
ओढली गेली नाही
तर सत्य असते;
पण ती ओढली
असती तर नक्कीच
तिचे फुलांच्या
गुच्छात रूपांतर
झाले नसते.

%READERNOTE{OLINC-238: खालील एका सर्वाधिक जवळच्या जगावर आधारित वर्णन हे दोन्ही सिद्धांतांचे सुलभीकरण आहे. स्टालनाकरचे विश्लेषण अद्वितीय जवळचे जग गृहीत धरते; लुईसच्या विश्लेषणात अनेक समसमान जवळची जगे असू शकतात, किंवा सर्वात जवळचे जग अस्तित्वातच नसेल. म्हणून हे वर्णन लुईसच्या पूर्ण सत्यता-अटी समजू नये.}
प्रतिवास्तविक विधानांचे
योग्य तर्कशास्त्र नेमके
काय, यावर अजूनही
वाद आहे. स्टालनाकर
आणि लुईस यांनी
एक प्रभावी विश्लेषण
दिले. त्यांच्या विचारांची
ढोबळ मांडणी अशी:
‘‘जर~$S$ असे घडले
असते, तर~$T$ असे
घडले असते’’ हे
प्रतिवास्तविक विधान
तेव्हा आणि तेव्हाच
सत्य असते जेव्हा
वास्तविक जगाशी
सर्वाधिक जवळीक
असलेल्या आणि~$S$
सत्य असलेल्या
प्रतिवास्तविक
परिस्थितीत
(‘‘संभाव्य जगात’’)
$T$ सत्य असते.
संभाव्य जगांच्या
सत्ताशास्त्राचा
संदर्भ घेत असल्याने
याला ‘‘सत्ताविषयक’’
विश्लेषण म्हणतात.
इतर विश्लेषणे सशर्त
संभाव्यता किंवा
विश्वास-पुनरीक्षणाचे
सिद्धांत वापरतात.
प्रतिवास्तविक विधानांसाठी
अनेक वेगवेगळी
तर्कशास्त्रे सुचवली
गेली आहेत. लुईस--स्टालनाकर
यांचे एकच तर्कशास्त्रही
नाही: सर्वाधिक जवळच्या
संभाव्य जगांच्या आधारे
त्यांनी साधर्म्य असलेली
मांडणी केली, तरी
त्यांच्या भिन्न
गृहीतकांमुळे भिन्न
वैध अनुमाने मिळतात.

\end{document}
